% ICASSP 2027: official spconf / IEEEbib; four content pages plus references.
\documentclass{article}
\usepackage{spconf,amsmath,amssymb,graphicx,hyperref}
\usepackage{booktabs,url,orcidlink}
\hypersetup{hidelinks}
% Generated by scripts/build_paper_numbers.py; do not edit display values.
\newcommand{\R}[1]{\csname result#1\endcsname}
\expandafter\def\csname resulttrainN\endcsname{102}
\expandafter\def\csname resulttrainPositive\endcsname{29}
\expandafter\def\csname resultdevN\endcsname{33}
\expandafter\def\csname resultdevPositive\endcsname{12}
\expandafter\def\csname resulttestN\endcsname{45}
\expandafter\def\csname resulttestPositive\endcsname{14}
\expandafter\def\csname resultcleanN\endcsname{180}
\expandafter\def\csname resultfitC\endcsname{0.3}
\expandafter\def\csname resultnestedN\endcsname{135}
\expandafter\def\csname resultall24Fmean\endcsname{0.635}
\expandafter\def\csname resultall24Amean\endcsname{0.676}
\expandafter\def\csname resultno_askFmean\endcsname{0.566}
\expandafter\def\csname resultno_askAmean\endcsname{0.610}
\expandafter\def\csname resultlatency_structureFmean\endcsname{0.550}
\expandafter\def\csname resultnestedRepeats\endcsname{10}
\expandafter\def\csname resultall24Dim\endcsname{24}
\expandafter\def\csname resultall24C\endcsname{0.30}
\expandafter\def\csname resultall24Dev\endcsname{0.746}
\expandafter\def\csname resultall24Test\endcsname{0.631}
\expandafter\def\csname resultno_askDim\endcsname{15}
\expandafter\def\csname resultno_askC\endcsname{1.00}
\expandafter\def\csname resultno_askDev\endcsname{0.814}
\expandafter\def\csname resultno_askTest\endcsname{0.641}
\expandafter\def\csname resultres_onlyDim\endcsname{8}
\expandafter\def\csname resultres_onlyC\endcsname{1.00}
\expandafter\def\csname resultres_onlyDev\endcsname{0.607}
\expandafter\def\csname resultres_onlyTest\endcsname{0.517}
\expandafter\def\csname resultask_onlyDim\endcsname{9}
\expandafter\def\csname resultask_onlyC\endcsname{0.01}
\expandafter\def\csname resultask_onlyDev\endcsname{0.594}
\expandafter\def\csname resultask_onlyTest\endcsname{0.593}
\expandafter\def\csname resultno_crossDim\endcsname{17}
\expandafter\def\csname resultno_crossC\endcsname{0.10}
\expandafter\def\csname resultno_crossDev\endcsname{0.700}
\expandafter\def\csname resultno_crossTest\endcsname{0.527}
\expandafter\def\csname resultablationdevpoint\endcsname{-0.068}
\expandafter\def\csname resultablationdevci_low\endcsname{-0.216}
\expandafter\def\csname resultablationdevci_high\endcsname{+0.083}
\expandafter\def\csname resultablationtestpoint\endcsname{-0.010}
\expandafter\def\csname resultablationtestci_low\endcsname{-0.089}
\expandafter\def\csname resultablationtestci_high\endcsname{+0.054}
\expandafter\def\csname resultask_dcoef\endcsname{-0.867}
\expandafter\def\csname resultask_dci_low\endcsname{-1.290}
\expandafter\def\csname resultask_dci_high\endcsname{-0.440}
\expandafter\def\csname resultres_hcoef\endcsname{+0.210}
\expandafter\def\csname resultres_hci_low\endcsname{-0.322}
\expandafter\def\csname resultres_hci_high\endcsname{+0.704}
\expandafter\def\csname resultAThreshold\endcsname{0.65}
\expandafter\def\csname resultADev\endcsname{0.673}
\expandafter\def\csname resultATestpoint\endcsname{0.545}
\expandafter\def\csname resultATestci_low\endcsname{0.392}
\expandafter\def\csname resultATestci_high\endcsname{0.681}
\expandafter\def\csname resultTThreshold\endcsname{0.5238}
\expandafter\def\csname resultTDev\endcsname{0.690}
\expandafter\def\csname resultTTestpoint\endcsname{0.631}
\expandafter\def\csname resultTTestci_low\endcsname{0.474}
\expandafter\def\csname resultTTestci_high\endcsname{0.769}
\expandafter\def\csname resultCTDThreshold\endcsname{0.50}
\expandafter\def\csname resultCTDDev\endcsname{0.746}
\expandafter\def\csname resultCTDTestpoint\endcsname{0.631}
\expandafter\def\csname resultCTDTestci_low\endcsname{0.472}
\expandafter\def\csname resultCTDTestci_high\endcsname{0.771}
\expandafter\def\csname resultTCTDThreshold\endcsname{0.55}
\expandafter\def\csname resultTCTDDev\endcsname{0.804}
\expandafter\def\csname resultTCTDTestpoint\endcsname{0.669}
\expandafter\def\csname resultTCTDTestci_low\endcsname{0.509}
\expandafter\def\csname resultTCTDTestci_high\endcsname{0.806}
\expandafter\def\csname resultfusionDelta\endcsname{0.038}
\expandafter\def\csname resultrobertaWeight\endcsname{0.3}
\expandafter\def\csname resultctdWeight\endcsname{0.7}
\expandafter\def\csname resultacousticWeight\endcsname{0.0}
\expandafter\def\csname resultmeandev\endcsname{0.738}
\expandafter\def\csname resultmeantest\endcsname{0.650}
\expandafter\def\csname resultmeanThreshold\endcsname{0.55}
\expandafter\def\csname resultpdchN\endcsname{62}
\expandafter\def\csname resultpdchSubjects\endcsname{46}
\expandafter\def\csname resultpdchRepeated\endcsname{16}
\expandafter\def\csname resultpdchPositive\endcsname{27}
\expandafter\def\csname resultpdchFmean\endcsname{0.520}
\expandafter\def\csname resultpdchFsd\endcsname{0.055}
\expandafter\def\csname resultpdchAmean\endcsname{0.476}
\expandafter\def\csname resultpdchAsd\endcsname{0.066}
\expandafter\def\csname resultpdchRepeats\endcsname{20}
\expandafter\def\csname resultdaicLatencyr\endcsname{0.245}
\expandafter\def\csname resultdaicLatencyq\endcsname{0.051}
\expandafter\def\csname resultpdchLatencyr\endcsname{0.267}
\expandafter\def\csname resultpdchLatencyq\endcsname{0.431}
\expandafter\def\csname resultbinaryHits\endcsname{0}
\expandafter\def\csname resultsignAgree\endcsname{12}
\expandafter\def\csname resultliveDim\endcsname{22}
\expandafter\def\csname resultmovingLatencyr\endcsname{0.107}
\expandafter\def\csname resultmovingLatencyq\endcsname{0.520}
\expandafter\def\csname resultmaxBinaryP\endcsname{0.340}
\expandafter\def\csname resultfdr\endcsname{0.10}
\expandafter\def\csname resultpermutations\endcsname{5000}
\expandafter\def\csname resultsilenceSeconds\endcsname{0.2}
\expandafter\def\csname resultrankLive\endcsname{13}
\expandafter\def\csname resultpowerN\endcsname{100}
\expandafter\def\csname resultpowerFmean\endcsname{0.605}
\expandafter\def\csname resultpowerFsd\endcsname{0.063}
\expandafter\def\csname resultpowerAmean\endcsname{0.636}
\expandafter\def\csname resultpowerAsd\endcsname{0.075}
\expandafter\def\csname resultpowerFpercentile\endcsname{10}
\expandafter\def\csname resultofficialtrainN\endcsname{107}
\expandafter\def\csname resultofficialdevN\endcsname{35}
\expandafter\def\csname resultofficialtestN\endcsname{47}
\expandafter\def\csname resultofficialN\endcsname{189}
\expandafter\def\csname resultexcludedN\endcsname{9}
\expandafter\def\csname resultnoPairsN\endcsname{3}
\expandafter\def\csname resultofficialC\endcsname{1.0}
\expandafter\def\csname resultofficialDev\endcsname{0.782}
\expandafter\def\csname resultofficialTest\endcsname{0.639}
\expandafter\def\csname resultofficialAuc\endcsname{0.654}
\expandafter\def\csname resultagarwalTest\endcsname{0.80}
\expandafter\def\csname resultwavlmSeedmean\endcsname{0.506}
\expandafter\def\csname resultwavlmSeedsd\endcsname{0.034}
\expandafter\def\csname resultwavlmSeeds\endcsname{6}
\expandafter\def\csname resultrobertaSeedmean\endcsname{0.604}
\expandafter\def\csname resultrobertaSeedsd\endcsname{0.025}
\expandafter\def\csname resultrobertaSeeds\endcsname{3}
\expandafter\def\csname resultwavlmSeed\endcsname{44}
\expandafter\def\csname resultouterFolds\endcsname{5}
\expandafter\def\csname resultinnerFolds\endcsname{4}
\expandafter\def\csname resultbootstrapN\endcsname{2000}
\expandafter\def\csname resultconfidence\endcsname{95}
\expandafter\def\csname resultthreshold\endcsname{0.5}
\expandafter\def\csname resultphqCutoff\endcsname{10}
\expandafter\def\csname resultcorrectedSession\endcsname{409}
\expandafter\def\csname resulthamdCutoff\endcsname{17}
\expandafter\def\csname resultannotationSeconds\endcsname{1}
\expandafter\def\csname resultcGrid\endcsname{0.01, 0.03, 0.1, 0.3, 1.0}
\expandafter\def\csname resultweightStep\endcsname{0.1}
\expandafter\def\csname resultthresholdLow\endcsname{0.05}
\expandafter\def\csname resultthresholdHigh\endcsname{0.95}
\expandafter\def\csname resultthresholdStep\endcsname{0.05}
\expandafter\def\csname resultunitSum\endcsname{1}
\expandafter\def\csname resultcrossDim\endcsname{7}

\title{Conversational Timing for Depression Screening: A Bounded Empirical Study}
\name{Hanie Kang\scalebox{0.75}{\orcidlink{0009-0005-8023-9672}}\textsuperscript{1},
Huang-Cheng Chou\scalebox{0.75}{\orcidlink{0000-0003-2125-5689}}\textsuperscript{2},
Sudarsana Reddy Kadiri\scalebox{0.75}{\orcidlink{0000-0001-5806-3053}}\textsuperscript{2},
Shrikanth Narayanan\scalebox{0.75}{\orcidlink{0000-0002-1052-6204}}\textsuperscript{2}}
\address{\textsuperscript{1}Thomas Lord Department of Computer Science, University of Southern California, USA\\
\textsuperscript{2}Signal Analysis and Interpretation Laboratory (SAIL), University of Southern California, USA}

\begin{document}
\maketitle
\begin{abstract}
We audit conversational timing (CTD) for depression screening using established descriptors. On DAIC-WOZ, nested macro-F1 is \R{all24Fmean} for full CTD, \R{latency_structureFmean} for latency plus interview-structure controls, and \R{no_askFmean} after interviewer-side removal. CTD carries information beyond these simple controls, with sensitivity to interviewer features and evaluation protocol. Recorded text--CTD fusion gains remain statistically inconclusive. On PDCH, latency--severity associations are weak and directionally consistent with DAIC-WOZ, but fail PDCH FDR; no binary feature passes FDR in both corpora, and sign agreement is \R{signAgree}/\R{liveDim}, near chance. Within-PDCH discrimination is unreliable (AUC \R{pdchAmean}). Language, instrument, population, interviewer, and timing measurement differ. This bounded empirical study does not claim detector transfer, construct validation, or a reliable multimodal advantage.
\end{abstract}
\begin{keywords}
depression screening, conversational timing, turn-taking, multimodal fusion, evaluation protocol
\end{keywords}

\section{Introduction}
Speech-based depression screening is studied as a low-cost complement to clinical assessment~\cite{valstar2016avec,cummins2015review}. Most systems model acoustic characteristics or lexical content. A clinical interview is also a dyadic interaction: an interviewer asks and a participant responds. Response delays, within-turn pauses, and turn durations describe when the exchange unfolds~\cite{yamamoto2020timing,aldeneh2019mood}.

Timing is jointly produced by both interlocutors. Interviewer prompts, session length, population, and annotation procedures can influence the same descriptors. A development-set gain therefore cannot by itself identify participant-intrinsic behavior or a clinically meaningful mechanism. Small interview corpora also make the apparent value of a feature sensitive to selection and evaluation choices.

We study conversational temporal dynamics (CTD), adopting \R{all24Dim} established turn-pair descriptors~\cite{chou2019joint,chou2021automatic}. Our bounded empirical contribution has three parts: (i) DAIC-WOZ timing information beyond simple latency and interview structure, with sensitivity to interviewer features and protocol; (ii) observed but statistically inconclusive fusion gains; and (iii) PDCH as a boundary, with unreliable within-corpus discrimination (AUC \R{pdchAmean}). Its weak, directionally consistent latency--severity association fails PDCH FDR; binary dual-corpus FDR hits are \R{binaryHits}, and signs agree for \R{signAgree}/\R{liveDim} features, near chance. Language, instrument, population, interviewer, and timing measurement differ. We claim neither detector transfer nor construct validation nor a reliable multimodal advantage. The repository documents the analysis pipeline.\footnote{\small\href{https://github.com/dndbsl/depression-detection-ctd}{\nolinkurl{github.com/dndbsl/depression-detection-ctd}}}

\section{Related Work}
Timing has substantial precedent: Aldeneh et al.~\cite{aldeneh2019mood} study dyadic turn-taking for mood episodes, and Fushimi et al.~\cite{fushimi2026beyond} study session-level utterance, pause, and response intervals for depression detection. Chou et al.~\cite{chou2019joint,chou2021automatic} introduce the turn-pair descriptors adopted here in deception detection. Our contribution concerns the scope of their empirical evidence, rather than feature novelty. Interviewer-prompt shortcuts on DAIC-WOZ~\cite{burdisso2024prompts} motivate explicit structure controls and cautious interpretation of dyadic features.

Frozen WavLM~\cite{chen2022wavlm} and RoBERTa~\cite{liu2019roberta} supply acoustic and text probes. Their results depend on pooling, head training, and selection. Agarwal et al.~\cite{agarwal2024discourse} provide published text-system context in Table~\ref{tab:context}; their use of both speakers' text and a different preprocessing protocol prevents a matched comparison. That method has not been rerun here.

\section{Data and Timing Representation}
\subsection{Corpora and integrity}
DAIC-WOZ~\cite{gratch2014daic} contains English interviews with Ellie, a human-controlled virtual interviewer, and self-reported PHQ-8 scores~\cite{kroenke2009phq8}. We predict PHQ-8 $\geq\R{phqCutoff}$, a screening label rather than a diagnosis. From the official \R{officialN} sessions, our integrity policy excludes \R{excludedN} sessions with timestamp problems, interruptions, or missing interviewer turns documented by the upstream preprocessing resource~\cite{baileyDaicProcessing}. Session \R{correctedSession} is retained with its binary label corrected to agree with its PHQ-8 score of \R{phqCutoff}. The cleaned cohort is \R{trainN}/\R{devN}/\R{testN} train/dev/test, with \R{trainPositive}/\R{devPositive}/\R{testPositive} positive labels; no subject changes official partition.

PDCH~\cite{cao2025pdch} comprises Mandarin clinician-led consultations with HAMD-17 ratings~\cite{hamilton1960rating}. We use \R{pdchN} labeled sessions from \R{pdchSubjects} subjects, including \R{pdchRepeated} subjects with repeated sessions. The study-specific target HAMD-17 $\geq\R{hamdCutoff}$ gives \R{pdchPositive} positives; continuous severity is analyzed separately. PHQ-8 and HAMD-17 are never pooled or treated as interchangeable labels. Chunk stitching and monaural voice-activity detection refine PDCH's \R{annotationSeconds}\,s turn annotations. Timing accuracy against manual speaker boundaries remains unvalidated. Overlap counts \texttt{ask\_bt} and \texttt{res\_bt} are structurally constant under this measurement pipeline, leaving \R{liveDim} measurable coordinates for the association audit. Language, instrument, population, interviewer, and timing measurement all differ from DAIC-WOZ.

\subsection{Fixed CTD descriptors}
Alternating interviewer (Ask) and participant (Res) runs form consecutive turn pairs. For each side, $d$ spans first onset to last offset, $u$ sums utterance durations, and $s=d-u$ describes internal gaps. Thus $u$ is an annotation-based speech-activity proxy. The descriptor bank includes both turn durations, their differences, sum and reciprocal ratios; all ordered quotients among $d,u,s$ for each side; response latency \texttt{res\_h} from Ask-end to Res-start; overlap counts; and within-turn gap counts, using \R{silenceSeconds}\,s on DAIC-WOZ. Undefined ratios remain missing until model fitting.

We average each descriptor over session turn pairs. Fixed groups contain \R{ask_onlyDim} ask-side, \R{crossDim} cross-speaker, and \R{res_onlyDim} response-internal coordinates. No-ask retains the cross-speaker group, which still uses interviewer timing. These coordinates are dependent: $u/d+s/d=\R{unitSum}$, for example. The centered matrices restricted to the \R{liveDim} live coordinates have observed rank \R{rankLive} on both corpora; we do not interpret coefficients as independent causal contributions.

\section{Evaluation Protocol}
\subsection{Original split and recorded neural runs}
CTD uses median imputation, standardization, and class-balanced $L_2$ logistic regression. On the original cleaned split, development balanced accuracy selects $C$ from the candidates \R{cGrid}, yielding $C=\R{fitC}$. The decision threshold is \R{threshold}. Dev predictions use a train-only fit and test predictions a train+dev refit, including preprocessing. Each feature-group ablation reselects $C$ under this same rule.

The acoustic probe pools frozen WavLM-large representations over participant utterances using a learned layer mixture and attention pooling. The text probe mean-pools frozen RoBERTa-large participant-turn embeddings with training-fitted standardization and a linear head. Interviewer lexical content is excluded from this text branch. Checkpoints were selected by dev loss for WavLM and dev macro-F1 for RoBERTa. However, WavLM seed \R{wavlmSeed} was chosen using \emph{test} macro-F1. Table~\ref{tab:results} therefore describes recorded runs and retains the \R{wavlmSeeds}-seed mean. Frozen caches, selected neural checkpoints, and session predictions are unavailable for a fresh rerun; no substitute seed is introduced. These are probe-specific results, not a ceiling for acoustic or text modeling.

Convex fusion selects nonnegative probability weights summing to one, on a dev grid of step \R{weightStep}. Thresholds are selected over \R{thresholdLow}--\R{thresholdHigh} in steps of \R{thresholdStep}, maximizing dev macro-F1. Equal-weight fusion also tunes its threshold. Component operating points are shown in Table~\ref{tab:results}; RoBERTa searches candidate probabilities, whereas CTD fixes its threshold. Recorded \R{confidence}\% intervals use \R{bootstrapN} participant bootstrap resamples~\cite{ferrer2024confidence}. They condition on fitted predictions and do not account for all training or selection uncertainty; dev intervals additionally reuse the selection set.

\subsection{Nested controls and association audit}
The exploratory timing-control protocol was specified after prior benchmark inspection. It uses only \R{nestedN} DAIC train+dev subjects, \R{outerFolds} outer folds, \R{innerFolds} inner folds, and \R{nestedRepeats} identical repeats across configurations. Inner balanced accuracy selects $C$; preprocessing is fitted within training folds and the decision threshold stays at \R{threshold}. Controls are latency alone, structure alone (log session span and log of one plus turn-pair count), and their combination. Out-of-fold predictions are pooled within each repeat. Repeats share subjects: their SDs are not confidence intervals or independent-repeat significance tests.

Within-PDCH prediction uses subject-grouped nested cross-validation, \R{outerFolds} outer and \R{innerFolds} inner folds over \R{pdchRepeats} repeats, selecting $C$ and the silence-gap threshold inside training folds. Both sessions of a subject stay together. Every model is fitted within its own corpus; this is not frozen detector transfer.

Association estimates use DAIC train+dev and the labeled PDCH cohort separately. The declared families cover binary labels, continuous severity, and psychomotor items, with Benjamini--Hochberg FDR at $q=\R{fdr}$ within family and corpus. PDCH uses \R{permutations} subject-block permutations. Prior exposure includes favorable DAIC latency test results and PDCH binary associations; later pre-declaration and nested fitting do not erase that exposure. Item exploration is not promoted as a finding.

\section{Results}
\subsection{Original-split prediction: fusion remains inconclusive}
Table~\ref{tab:results} preserves the recorded modality results while the canonical CTD baseline is rerun. Text--CTD fusion has test macro-F1 \R{TCTDTestpoint}, an observed increase of \R{fusionDelta} over either component's rounded \R{CTDTestpoint}. Recorded paired intervals include or touch zero; the missing session predictions prevent refreshing their exact limits. Thus the fusion gain remains statistically inconclusive. Equal-weight fusion gives \R{meandev}/\R{meantest} dev/test, with its threshold still selected on dev. The text probe's recorded test mean is $\R{robertaSeedmean}\pm\R{robertaSeedsd}$ over \R{robertaSeeds} seeds.

\begin{table}[t]
\centering
\caption{Cleaned DAIC-WOZ macro-F1. Neural/fusion rows describe recorded runs; CTD was reproduced. Dev is used for selection. Test brackets are \R{confidence}\% bootstrap CIs, except the final seed mean$\pm$SD. Thresholds $\tau$ are shown rounded. A: WavLM; T: RoBERTa.}
\label{tab:results}
\small
\setlength{\tabcolsep}{3pt}
\begin{tabular}{lccc}
\toprule
System & $\tau$ & Dev & Test [CI]\\
\midrule
A (seed \R{wavlmSeed}) & \R{AThreshold} & \R{ADev} & \R{ATestpoint} [\R{ATestci_low},\R{ATestci_high}]\\
T & \R{TThreshold} & \R{TDev} & \R{TTestpoint} [\R{TTestci_low},\R{TTestci_high}]\\
CTD & \R{CTDThreshold} & \R{CTDDev} & \R{CTDTestpoint} [\R{CTDTestci_low},\R{CTDTestci_high}]\\
A+T & \R{ATThreshold} & \R{ATDev} & \R{ATTestpoint} [\R{ATTestci_low},\R{ATTestci_high}]\\
A+CTD & \R{ACTDThreshold} & \R{ACTDDev} & \R{ACTDTestpoint} [\R{ACTDTestci_low},\R{ACTDTestci_high}]\\
T+CTD & \R{TCTDThreshold} & \R{TCTDDev} & \R{TCTDTestpoint} [\R{TCTDTestci_low},\R{TCTDTestci_high}]\\
Mean(T+CTD) & \R{meanThreshold} & \R{meandev} & \R{meantest}\\
\midrule
A: \R{wavlmSeeds}-seed mean & -- & -- & $\R{wavlmSeedmean}\pm\R{wavlmSeedsd}$\\
\bottomrule
\end{tabular}
\end{table}

Three-way convex fusion selects weights A=\R{acousticWeight}, T=\R{robertaWeight}, and CTD=\R{ctdWeight}. It coincides with T+CTD, so we omit the duplicate row. Zero acoustic weight is specific to the probe, scores, threshold, and grid; it does not imply acoustic redundancy. A stronger acoustic rerun is outside this pass because its artifacts are unavailable.

The official \R{officialtrainN}/\R{officialdevN}/\R{officialtestN} sensitivity run restores all \R{excludedN} exclusions, retains the session-\R{correctedSession} label correction, and selects $C=\R{officialC}$. CTD gives dev/test macro-F1 \R{officialDev}/\R{officialTest} and test AUC \R{officialAuc}. The \R{noPairsN} sessions without Ask/Res pairs receive entirely missing descriptors and training-fitted median imputation. Timestamp and interruption defects remain. This check exposes sensitivity to the integrity policy; it does not certify timing quality. T and T+CTD cannot be evaluated on this cohort without the missing neural predictions.

\begin{table}[t]
\centering
\caption{Official-cohort context. These are \emph{not matched-protocol comparisons}: Agarwal et al. use both speakers' text; CTD uses timing and restores defective sessions as described in the text. No published method was reproduced.}
\label{tab:context}
\small
\begin{tabular}{lcc}
\toprule
System & Train/dev/test & Test macro-F1\\
\midrule
Agarwal et al.~\cite{agarwal2024discourse} & \R{officialtrainN}/\R{officialdevN}/\R{officialtestN} & \R{agarwalTest}\\
CTD sensitivity & \R{officialtrainN}/\R{officialdevN}/\R{officialtestN} & \R{officialTest}\\
\bottomrule
\end{tabular}
\end{table}

\subsection{Nested timing controls and interviewer sensitivity}
Full CTD has higher mean scores than latency plus structure (macro-F1 \R{all24Fmean} versus \R{latency_structureFmean}; Table~\ref{tab:controls}). Its information on this DAIC cohort is not exhausted by these simple controls. Removing the ask-side group lowers mean macro-F1 to \R{no_askFmean} and AUC from \R{all24Amean} to \R{no_askAmean}. These descriptive contrasts show sensitivity to ask-side removal; dependent repeats do not supply a significance test.

\begin{table}[t]
\centering
\caption{DAIC train+dev nested controls, all declared configurations. Mean$\pm$SD across \R{nestedRepeats} paired repeats; SD is not a CI. Official test subjects are excluded.}
\label{tab:controls}
\small
\begin{tabular}{lcc}
\toprule
Representation & Macro-F1 & ROC-AUC\\
\midrule
Full CTD & $\R{all24Fmean}\pm\R{all24Fsd}$ & $\R{all24Amean}\pm\R{all24Asd}$\\
No ask-side group & $\R{no_askFmean}\pm\R{no_askFsd}$ & $\R{no_askAmean}\pm\R{no_askAsd}$\\
Latency & $\R{latencyFmean}\pm\R{latencyFsd}$ & $\R{latencyAmean}\pm\R{latencyAsd}$\\
Structure & $\R{structureFmean}\pm\R{structureFsd}$ & $\R{structureAmean}\pm\R{structureAsd}$\\
Latency + structure & $\R{latency_structureFmean}\pm\R{latency_structureFsd}$ & $\R{latency_structureAmean}\pm\R{latency_structureAsd}$\\
\bottomrule
\end{tabular}
\end{table}

The original-split ablation tells a different descriptive story (Table~\ref{tab:ablation}). With $C$ reselected, no-ask gives \R{no_askDev}/\R{no_askTest}. The all-CTD minus no-ask macro-F1 difference is \R{ablationdevpoint} [\R{ablationdevci_low},\R{ablationdevci_high}] on dev and \R{ablationtestpoint} [\R{ablationtestci_low},\R{ablationtestci_high}] on test (paired \R{confidence}\% CIs). No clear difference is detected on that split; nonsignificance does not establish equivalence. Cross-speaker coordinates still encode interviewer timing, so the ablation cannot remove interviewer confounding.

\begin{table}[t]
\centering
\caption{Original cleaned split: feature-group ablations with $C$ reselected by dev balanced accuracy for every row. Scores are macro-F1; test comparisons revisit the same cohort.}
\label{tab:ablation}
\small
\setlength{\tabcolsep}{4pt}
\begin{tabular}{lrrrr}
\toprule
Configuration & Dim. & $C$ & Dev & Test\\
\midrule
Full CTD & \R{all24Dim} & \R{all24C} & \R{all24Dev} & \R{all24Test}\\
No ask-side & \R{no_askDim} & \R{no_askC} & \R{no_askDev} & \R{no_askTest}\\
No cross-speaker & \R{no_crossDim} & \R{no_crossC} & \R{no_crossDev} & \R{no_crossTest}\\
Response-internal & \R{res_onlyDim} & \R{res_onlyC} & \R{res_onlyDev} & \R{res_onlyTest}\\
Ask-side & \R{ask_onlyDim} & \R{ask_onlyC} & \R{ask_onlyDev} & \R{ask_onlyTest}\\
\bottomrule
\end{tabular}
\end{table}

Interviewer duration \texttt{ask\_d} has the largest train-fitted standardized coefficient: \R{ask_dcoef}, bootstrap CI [\R{ask_dci_low},\R{ask_dci_high}]. Latency's coefficient is \R{res_hcoef} [\R{res_hci_low},\R{res_hci_high}]. Correlated coordinates and protocol-sensitive ablations limit any claim that one feature identifies the mechanism of the multivariate detector.

\subsection{PDCH as a boundary}
Within-PDCH full CTD gives macro-F1 $\R{pdchFmean}\pm\R{pdchFsd}$ and AUC $\R{pdchAmean}\pm\R{pdchAsd}$, showing no reliable discrimination. Latency--severity Spearman correlations are weak and directionally consistent: DAIC $\rho=\R{daicLatencyr}$ ($q=\R{daicLatencyq}$), PDCH $\rho=\R{pdchLatencyr}$ ($q=\R{pdchLatencyq}$, failing FDR). In the binary family, \R{binaryHits} features pass FDR in both corpora; signs agree for \R{signAgree}/\R{liveDim}, near chance descriptively, and PDCH's max-$|r|$ permutation test gives $p=\R{maxBinaryP}$. Language, instrument, population, interviewer, and timing measurement differ, preventing attribution of this boundary to one factor.

For latency, the DAIC \texttt{PHQ8\_Moving} item is also inconclusive ($\rho=\R{movingLatencyr}$, $q=\R{movingLatencyq}$). This self-report item combines slowing and restlessness; its null does not isolate a psychomotor mechanism or establish inferiority to clinician ratings. The PDCH item exploration is not treated as validation. Small-sample DAIC downsampling, matching, and residualization are retained in the audit; they do not identify low power as the sole explanation for PDCH or rule out other confounding.

\section{Discussion and Limitations}
The bounded conclusion is that CTD carries DAIC-WOZ information beyond simple latency and structure controls, with interviewer-feature and protocol sensitivity; observed fusion gains remain inconclusive. PDCH supplies a limit: weak directional latency--severity concordance fails its FDR threshold, with \R{binaryHits} dual-corpus binary hits, \R{signAgree}/\R{liveDim} sign agreement near chance, and unreliable discrimination (AUC \R{pdchAmean}). Language, instrument, population, interviewer, and timing measurement change together. These results provide neither detector transfer nor construct validation nor a reliable multimodal advantage.

Timing-only nested controls do not assess fusion stability. The declared text, text+CTD, text+no-ask, and text+simple-controls comparison on the same outer folds remains unrun because raw frozen embeddings are missing. A stronger text/acoustic baseline and independent, label-blind timing-boundary annotation would address important remaining limits. Fresh disjoint participants are needed for external evaluation; E-DAIC is a possible future resource, requiring speaker labeling and overlap checks, not a completed validation here. Subgroup performance is unmeasured. These screening-label models are not validated for clinical deployment; consent, privacy, and demographic performance require separate assessment.

\section{Conclusion}
CTD carries information on DAIC-WOZ beyond simple latency and interview-structure controls, but is sensitive to interviewer features and evaluation protocol. Fusion gains are observed but statistically inconclusive. PDCH's weak, directionally consistent latency--severity association fails FDR, alongside \R{binaryHits} dual-corpus binary hits, \R{signAgree}/\R{liveDim} sign agreement near chance, and no reliable within-corpus discrimination (AUC \R{pdchAmean}); language, instrument, population, interviewer, and timing measurement differ. We do not claim detector transfer, construct validation, or a reliable multimodal advantage.

\section{Compliance with Ethical Standards}
This is secondary analysis of existing DAIC-WOZ and PDCH interviews, with no new participant recruitment. Original data collection and consent procedures are described by the corpus authors~\cite{gratch2014daic,cao2025pdch}.
% AUTHOR CHECK BEFORE SUBMISSION: confirm institutional approval/waiver or
% data-use basis, plus funding and conflicts. No unknown approval is asserted.

\bibliographystyle{IEEEbib}
\bibliography{references}
\end{document}
